MHC-I presentation can vary among tissues even when the presenting allele and source protein are fixed. Existing predictors largely estimate tissue-blind peptide--MHC compatibility and therefore do not directly address this tissue-associated preference. We formulate a matched-pair ranking problem: given two peptides from the same parent protein under the same presenting restriction, determine which has stronger presentation evidence in a specified tissue. The benchmark contains 157 human tissue--HLA combinations and 24 mouse tissue--H-2 combinations, where H-2 is the mouse major histocompatibility complex; each contains more than 200 matched peptide pairs. TissuePMHC combines a peptide representation learned across all human tasks with an allele-restricted representation, while species-specific baselines and dual-branch models are evaluated independently in mouse. On saved internal test sets, the strongest retained prediction aggregates have equal-weight mean AUROC/AUPRC values of 0.8448/0.8348 in human and 0.8528/0.8477 in mouse. When every held-out peptide identity is excluded from all training tasks, the highest identical-fold AUROC values are 0.7652 and 0.7590, respectively. Removing tissue identity substantially reduces human performance, showing that general peptide--MHC compatibility alone does not reproduce the complete ranking signal. The full human architecture is statistically comparable to its strongest simpler guided control, while mouse auxiliary supervision provides a small but consistent improvement over the matched unguided dual branch. These results support learnability of tissue-conditioned peptide ranking for previously represented tissue--MHC combinations and quantify the optimism associated with cross-task peptide reuse in human. They do not establish a causal tissue-processing mechanism, generalization to unseen tissues or alleles, independent-cohort validity, or clinical utility.
